\documentclass{article}
\usepackage[UTF8]{ctex}
\usepackage{amsmath,amssymb,amsfonts}
\usepackage{enumitem}
\usepackage{geometry}
\geometry{a4paper, margin=1in}
% === 错题本排版优化（不参与题目标识） ===
\setlength{\parindent}{0pt}
\setlength{\parskip}{0.35em plus 0.10em minus 0.08em}
\setlength{\emergencystretch}{3em}
\setlist{itemsep=0.25em, topsep=0.40em, parsep=0pt}
\allowdisplaybreaks[2]
% === 错题本排版优化结束 ===
\begin{document}
\section*{武汉理工大学2025年数学分析真题}
\section*{136．武汉理工大学2025年数学分析试卷}

\subsection{题目1}
1．计算极限 \(\lim _{x \rightarrow 0} \frac{\sqrt[3]{1+\tan ^{2} x}-\sqrt{\cos x}}{x \sin x}\) ．

\subsection{题目2}
2．设 \(f(u, v)\) 在 \(\mathbb{R}^{2}\) 上有二阶连续偏导数，且 \(\frac{\partial^{2} f}{\partial u^{2}}+\frac{\partial^{2} f}{\partial v^{2}}=1\) ，设 \(g(x, y)=f\left(x y, \frac{x^{2}-y^{2}}{2}\right)\) ，求

\[
\frac{\partial^{2} g}{\partial x^{2}}+\frac{\partial^{2} g}{\partial y^{2}} .
\]

\subsection{题目3}
3．计算 \(\iint_{D} \frac{x \sin ^{2} y}{y^{2}} \mathrm{~d} x \mathrm{~d} y\) ，其中 \(D\) 是由 \(y=x, x=0, y=\pi\) 围成的区域。

\subsection{题目4}
4．计算 \(\int_{L} y e^{y^{2}} \mathrm{~d} y+\left(x e^{y^{2}}+2 x y^{2} e^{y^{2}}\right) \mathrm{d} x\) ，其中 \(L: y=\sqrt[3]{x}\) 从 \((0,0)\) 到 \((1,1)\) ．

\subsection{题目5}
5．设曲线 \(C:\left\{\begin{array}{l}x^{2}+y^{2}=2 a z ; \\ x^{2}+y^{2}+x y=a^{2} .\end{array}\right.\) 求曲线上点到坐标系平面 \(x O y\) 的最大值与最小值．

\subsection{题目6}
6．讨论级数 \(\sum_{n=1}^{\infty} \frac{2^{n}}{n} x^{2 n+1}\) 的收敛域并求其和函数。

\subsection{题目7}
7．设 \(\varphi(x, y)\) 在 \((0,0)\) 的某邻域内连续，\(f(x, y)=|x-y| \varphi(x, y)\) ．

\begin{itemize}
  \item[（1）] 讨论 \(\varphi(x, y)\) 满足什么条件时，\(f_{x}(0,0), f_{y}(0,0)\) 存在．
  \item[（2）] 讨论 \(\varphi(x, y)\) 满足什么条件时，\(f(x, y)\) 在 \((0,0)\) 处可微．
\end{itemize}

\subsection{题目8}
8．讨论反常积分 \(\int_{1}^{+\infty} \frac{\sin x}{x^{p}} \mathrm{~d} x(p>0)\) 的敛散性，并说明收敛时是绝对收敛还是条件收敛．

\subsection{题目9}
9．设 \(f(x)\) 在 \((-\infty,+\infty)\) 上一致连续，\(f(x) \in[a, b], g(x)\) 在 \([a, b]\) 上连续，证明：\(g(f(x))\) 在\((-\infty,+\infty)\) 上一致连续．

\subsection{题目10}
10．证明：\(f(x)\) 在 \([a, b]\) 上连续，则 \(f(x)\) 在 \([a, b]\) 上有界．

\subsection{题目11}
11．设数列 \(\left\{a_{n}\right\}\) 为恒正的单调增数列，证明：

\begin{itemize}
  \item[（1）] 若 \(\left\{a_{n}\right\}\) 有界，则 \(\sum_{n=1}^{\infty}\left(1-\frac{a_{n}}{a_{n+1}}\right)\) 收敛。
  \item[（2）] 若 \(\left\{a_{n}\right\}\) 无界，则 \(\sum_{n=1}^{\infty}\left(1-\frac{a_{n}}{a_{n+1}}\right)\) 发散。
\end{itemize}

\subsection{题目12}
12．设 \(f(x)\) 在 \([0,+\infty)\) 上连续，\(A>0, \int_{A}^{+\infty} \frac{f(x)}{x} \mathrm{~d} x\) 存在，证明：对任意的 \(a, b>0\) ，有

\[
\int_{0}^{+\infty} \frac{f(b x)-f(a x)}{x} \mathrm{~d} x=f(0) \ln \frac{a}{b} .
\]
\end{document}
